\documentclass{article}

% if you need to pass options to natbib, use, e.g.:
%     \PassOptionsToPackage{numbers, compress}{natbib}
% before loading neurips_2024


% ready for submission
\usepackage[preprint,nonatbib]{neurips_2024}


% to compile a camera-ready version, add the [final] option, e.g.:
%     \usepackage[final]{neurips_2024}

\usepackage[utf8]{inputenc} % allow utf-8 input
\usepackage[T1]{fontenc}    % use 8-bit T1 fonts
\usepackage{hyperref}       % hyperlinks
\usepackage{url}            % simple URL typesetting
\usepackage{booktabs}       % professional-quality tables
\usepackage{amsfonts}       % blackboard math symbols
\usepackage{nicefrac}       % compact symbols for 1/2, etc.
\usepackage{microtype}      % microtypography
\usepackage{xcolor}         % colors
\usepackage{todonotes}

\usepackage[backend=biber,style=numeric]{biblatex}

\addbibresource{Masterseminar.bib}

\title{Evaluating the Decision-Making Capabilities of LLM-based Agents for Production Scheduling}


% The \author macro works with any number of authors. There are two commands
% used to separate the names and addresses of multiple authors: \And and \AND.
%
% Using \And between authors leaves it to LaTeX to determine where to break the
% lines. Using \AND forces a line break at that point. So, if LaTeX puts 3 of 4
% authors names on the first line and the last on the second line, try using
% \AND instead of \And before the third author name.


\author{%
  Name Surename \\
  Informatik, Matr NR \\
  \texttt{emailaddress@tha.de} \\
}


\begin{document}


\maketitle

\section{Introduction and Motivation}
Industrial production systems are often characterized by a high level of complexity in production planning due to fluctuating workloads, variable processing times and limited machine availability.
Efficient production planning therefore represents a central challenge in manufacturing environments. In particular, flexible job-shop scheduling problems require the selection of machines and processing sequences for multiple jobs while taking resource constraints into account.

Traditional scheduling approaches are often based on predefined heuristics or rule-based strategies, such as priority rules (e.g., SPT or EDD)~\cite[Chapter 3]{leungHandbookSchedulingAlgorithms2004}. In these approaches, decisions are made according to fixed priority rules without adapting to the current system state. While these methods achieve good results in stable and deterministic environments, they reach their limits in dynamic production systems because they cannot flexibly respond to changes such as machine failures or new incoming jobs.

Advances in artificial intelligence have introduced new approaches for solving complex decision-making problems. Reinforcement learning (RL) has shown promising results for dynamic scheduling tasks~\cite{waschneckOptimizationGlobalProduction2023, waschneckDeepReinforcementLearning2018, zhangDeepReinforcementLearning2023}; however, RL-based solutions typically require extensive training and often need to be retrained when the production environment changes.

Large language models (LLMs) have recently demonstrated strong capabilities in reasoning, planning and interaction with tools~\cite{yaoREACSYNERGIZINGREASONING2023, schickToolformerLanguageModels}, particularly when deployed as agent-based systems~\cite{vyasAutonomousIndustrialControl2024a, xiaControlIndustrialAutomation2025a}. These models can process structured inputs, generate decisions and interact with external systems. LLM-based agents could therefore serve as flexible decision-making components in production control systems.

This work focuses on the systematic evaluation of LLM-based agents for production planning in flexible job-shop environments, with the aim of investigating their ability to generate valid and efficient scheduling decisions under different system configurations.



\section{State of the Art}

Classical approaches include rule-based heuristics and optimization methods that aim to minimize performance indicators such as makespan, waiting times, or machine idle times~\cite[Chapter 3]{leungHandbookSchedulingAlgorithms2004}.

More recently, reinforcement learning has been used to learn scheduling strategies through interaction with production environments. Zhang et al.~\cite{zhangDeepReinforcementLearning2023} investigate deep reinforcement learning for dynamic flexible job-shop scheduling with variable processing times.

Further work demonstrates the applicability of reinforcement learning in complex industrial scenarios, such as semiconductor manufacturing, where scheduling decisions must be made under high system complexity and uncertainty~\cite{waschneckDeepReinforcementLearning2018}.

In addition, reinforcement learning has also been applied to the optimization of global production planning problems, where multiple sites and complex dependencies must be taken into account~\cite{waschneckOptimizationGlobalProduction2023}.

Overall, these approaches show that reinforcement learning is capable of learning adaptive and high-performing scheduling strategies in dynamic and complex production environments.

At the same time, current research explores the use of large language models as agents that are capable of reasoning and interacting with external systems. Foundational work shows that large language models can combine reasoning and action capabilities and integrate external tools into decision-making processes. Yao et al.~\cite{yaoREACSYNERGIZINGREASONING2023}, for example, demonstrate that language models can structure and solve complex tasks by combining reasoning and acting. Schick et al.~\cite{schickToolformerLanguageModels} further show that language models can learn to use external tools autonomously in order to extend their capabilities.

Building on these capabilities, more recent studies investigate the use of LLM-based agents in industrial contexts. Xia et al.~\cite{xiaControlIndustrialAutomation2025b} propose a framework in which an LLM agent interacts with industrial automation systems via a tool-based architecture. Similarly, Vyas et al.~\cite{vyasAutonomousIndustrialControl2024a} present an agent-based framework for autonomous industrial control with LLMs.

These developments highlight the potential of LLM-based agents as flexible and adaptive decision-making components that are particularly suitable for complex and dynamic industrial applications.

Despite this progress, the application of LLM-based agents to production scheduling remains largely unexplored. In particular, empirical studies on the effectiveness, robustness and transferability of LLM-generated scheduling decisions in complex production environments are still lacking.



\section{Problem Definition and Research Questions}

Based on the current state of research, several open questions remain regarding the role of large language models in scheduling tasks.

This work therefore focuses on evaluating the decision-making capabilities of LLM-based agents in flexible job-shop environments. The following research questions guide the investigation:


\begin{itemize}

\item \textbf{Integration of LLM-based Agents}
\begin{itemize}
    \item Which integration architectures are the most promising for incorporating LLM-based agents into production planning systems?
\end{itemize}

\item \textbf{Performance and Real-Time Suitability}
\begin{itemize}
    \item To what extent can LLM-based agents generate valid and effective scheduling decisions in flexible job-shop environments?
    
    \item How suitable is the use of LLM-based agents for real-time planning and control of production systems, particularly with regard to response times and decision latency?
\end{itemize}

\item \textbf{Types of Scheduling Decisions}
\begin{itemize}
    \item Which types of scheduling decisions (e.g., job selection, machine assignment and sequencing) can be carried out reliably by LLM-based agents?
\end{itemize}

\item \textbf{Factors Influencing Decision Quality}
\begin{itemize}
    \item How does the representation of the production state (e.g., JSON format, tabular representation, or textual description) influence the quality of the generated scheduling decisions?
    
    \item How does increasing production system complexity affect the decision quality of LLM-based scheduling agents?
\end{itemize}

\item \textbf{Influence of Prompt Design}
\begin{itemize}
    \item How does prompt design affect the quality and consistency of the generated scheduling decisions?
\end{itemize}

\item \textbf{Comparison with Classical Methods (optional)}
\begin{itemize}
    \item (Optional) How does the performance of LLM-based scheduling approaches compare to classical heuristic methods?
\end{itemize}

\end{itemize}



\section{Planned Methodology}

The aim of this work is to analyze the performance of large language models as decision-making components in scheduling tasks. The study combines a literature review with an experimental evaluation in a simulated production environment.

A flexible job-shop simulation environment models a simplified industrial production system with multiple machines, different job types and dynamic job arrivals. Within this environment, a LLM-based scheduling agent observes the current system state and generates scheduling decisions, such as selecting the next job and assigning it to an available machine.

The LLM acts as a higher-level decision-making component within the scheduling process. It receives structured information about the production state and derives scheduling actions from it. The experiments systematically vary the input representation (e.g., JSON, tabular, or textual), the prompt design and the complexity of the production environment (e.g., number of machines, jobs and dynamic disruptions) in order to examine their influence on decision quality. Multiple runs are conducted for each configuration to analyze the stability and reproducibility of the generated decisions.

The evaluation is based on established scheduling metrics such as makespan, flow time and machine utilization. In addition, the validity and consistency of the generated decisions are analyzed, particularly with respect to compliance with system constraints and the reproducibility of decisions under identical states.

(Optional) The performance of the LLM-based scheduler is compared with classical heuristic scheduling strategies such as First Come First Serve (FCFS), Shortest Processing Time (SPT) and Earliest Due Date (EDD).

In addition, different types of scheduling decisions are examined, including job selection, machine assignment and integrated scheduling decisions, in order to evaluate the capabilities of the LLM across different decision dimensions.


\subsection{Data}

The experiments use simulated production data based on a flexible job-shop scheduling environment and inspired by a realistic production system with multiple machines and connected conveyor belts. The structure of the system is modeled as a network of machines (nodes) and transport connections (edges).

Synthetic datasets are generated to simulate dynamic production conditions, including varying job arrival rates, stochastic processing times, machine availability and possible disruptions such as machine failures.

The simulation environment provides the scheduling agent with the current system state, including information about machine states, waiting jobs, processing times, queues and transport times. This production state is provided to the LLM in different representations, for example as JSON structures, tabular formats, or natural language descriptions.

\subsection{Planned Model / Experiment / Analysis}

The experimental setup consists of a simulated flexible job-shop environment in which scheduling decisions must be made dynamically.

The LLM acts as a decision-making component that generates scheduling actions based on the current system state. These actions include job selection, assignment to a machine and the determination of processing order. The generated actions are then checked for validity and executed in the simulation system.

Production planning comprises several decision types, including job selection, machine assignment and integrated scheduling decisions, which are systematically investigated in the experiments.

The experiments systematically vary the input representation (e.g., JSON, tabular, textual), the prompt design and the complexity of the production environment (e.g., number of machines, jobs and disruptions). In addition, multiple runs are carried out per scenario to ensure the stability and reproducibility of the results.

The analysis focuses on the quantitative evaluation of scheduling performance as well as on the investigation of robustness, decision validity and sensitivity to different input representations and system configurations.

\subsection{Planned Evaluation}

The evaluation is based on common performance metrics from the production scheduling literature. These allow for a quantitative comparison between the LLM-based scheduling agent and baseline methods.

Possible evaluation metrics include:

\begin{itemize}
\item makespan (total completion time)
\item average flow time of jobs
\item machine utilization
\item job waiting time
\item number of tardy jobs
\end{itemize}

In addition to the quantitative metrics, LLM-specific evaluation criteria are considered. These include in particular the validity of the generated decisions (compliance with system constraints), the consistency of decisions under identical inputs and the sensitivity to different input representations and prompt variants.

A qualitative analysis of selected decision situations is also conducted in order to better understand the decision behavior of LLM-based agents in complex scheduling scenarios.

In addition, the decision time of the LLM is analyzed to assess its suitability for real-time applications.

\section{(Optional) Work Plan and Timeline}

\begin{itemize}
\item Weeks 1--2: literature review on production planning and LLM-based agents
\item Weeks 3--4: implementation and refinement of the LLM-based scheduling prototype
\item Weeks 5--6: experimental evaluation and comparison with heuristic baselines
\item Weeks 7--8: writing and revision of the seminar paper
\end{itemize}

\printbibliography


\end{document}